%! Solving Exponential Equations with a Common Base — Unit 6, Lesson 6.3 — Homework
\documentclass[10pt]{article}
\usepackage{algebra2-article}
\usepackage{algebra2-boxes}
\newcommand{\TallMath}[1]{$\displaystyle #1\rule[-1.4em]{0pt}{3.2em}$}
\newcommand{\lgrid}[4]{%
  \draw[step=1, linegray!40, very thin] (#1,#3) grid (#2,#4);
  \draw[->, semithick, charcoal] (#1,0)--(#2,0) node[below right, font=\scriptsize]{$x$};
  \draw[->, semithick, charcoal] (0,#3)--(0,#4) node[left, font=\scriptsize]{$y$};}
% #1=a  #2=ln(b)  #3=h  #4=k  #5=xmin  #6=xmax   (ln2 = 0.693147)
\newcommand{\expgraph}[6]{\draw[very thick, forest, domain=#5:#6, samples=140, smooth]
  plot (\x,{(#1)*exp((#2)*((\x)-(#3)))+(#4)});}
\newcommand{\solline}[3]{\draw[goldacc, very thick] (#1,#3)--(#2,#3);}

\begin{document}

\pageheader{Unit 6, Lesson 6.3}{Homework}
\namedateperiod

\vspace{0.10in}

\begin{notesbox}{Practice}
Every solution below needs a visible common-base step. An answer with no rewriting shown is not a
solution.

\begin{enumerate}[leftmargin=1.9em, itemsep=12pt]
  \item \textbf{Solve each equation.}
    \begin{center}
    \renewcommand{\arraystretch}{1.9}
    \begin{tabularx}{\linewidth}{@{}l X c@{\hspace{1.6em}} l X c@{}}
    \hline
    \textbf{(a)} $2^{x}=128$ & \blank{2.4cm} & $x=$ \blank{1.3cm} &
    \textbf{(b)} $3^{\,2x}=81$ & \blank{2.4cm} & $x=$ \blank{1.3cm} \\
    \textbf{(c)} $5^{x}=\frac{1}{25}$ & \blank{2.4cm} & $x=$ \blank{1.3cm} &
    \textbf{(d)} $16^{x}=64$ & \blank{2.4cm} & $x=$ \blank{1.3cm} \\
    \textbf{(e)} $\left(\frac12\right)^{x}=32$ & \blank{2.4cm} & $x=$ \blank{1.3cm} &
    \textbf{(f)} $27^{\,x-1}=9$ & \blank{2.4cm} & $x=$ \blank{1.3cm} \\
    \textbf{(g)} $4^{\,x+2}=8^{x}$ & \blank{2.4cm} & $x=$ \blank{1.3cm} &
    \textbf{(h)} $6^{\,3x}=\sqrt{6}$ & \blank{2.4cm} & $x=$ \blank{1.3cm} \\
    \hline
    \end{tabularx}
    \end{center}
    \emph{Work space --- show the rewriting for at least four of the eight.}
    \vspace{1.05in}

    Three of your eight answers are not whole numbers. List them, and say in one sentence why that is
    not a sign of a mistake. \blank{7.0cm}

  \item \textbf{Read it, then prove it.} The curve is $y=2^{x}$.
    \begin{center}
    \begin{minipage}[c]{0.40\linewidth}
    \centering
    \begin{tikzpicture}[scale=0.38]
      \lgrid{-3}{4}{-5}{9}
      \begin{scope}
        \clip (-3,-5) rectangle (4,9);
        \expgraph{1}{0.693147}{0}{0}{-3}{4}
      \end{scope}
      \solline{-3}{4}{8}
      \solline{-3}{4}{2}
      \draw[redacc, dashed, very thick] (-3,-4)--(4,-4);
      \node[goldacc, font=\scriptsize, left] at (-1.2,8.6) {$y=8$};
      \node[goldacc, font=\scriptsize, left] at (-1.2,2.6) {$y=2$};
      \node[redacc, font=\scriptsize, left] at (-1.2,-4.6) {$y=-4$};
    \end{tikzpicture}
    \end{minipage}
    \begin{minipage}[c]{0.56\linewidth}
    \small
    \renewcommand{\arraystretch}{1.7}
    \begin{tabularx}{\linewidth}{@{}l X c@{}}
    \hline
    \textbf{Equation} & \textbf{Common base} & \textbf{$x=$} \\
    \hline
    $2^{x}=8$ & \blank{2.0cm} & \blank{1.2cm} \\
    $2^{x}=2$ & \blank{2.0cm} & \blank{1.2cm} \\
    $2^{x}=-4$ & \blank{2.0cm} & \blank{1.2cm} \\
    \hline
    \end{tabularx}
    \end{minipage}
    \end{center}
    \vspace{2pt}
    \textbf{(a)} What horizontal line would you have to draw to make the solution $x=0$?
    \blank{3.0cm} \\[4pt]
    \textbf{(b)} What horizontal line would give the solution $x=-1$? \blank{3.0cm} \\[4pt]
    \textbf{(c)} Explain why \emph{no} horizontal line drawn below the $x$-axis will ever produce a
    solution. \blank{6.5cm}

  \item \textbf{Sort them.} Decide whether each equation can be solved by writing both sides over a
    common base. Write \textbf{yes} or \textbf{no}, and solve the ones you can.
    \begin{center}
    \renewcommand{\arraystretch}{1.75}
    \begin{tabularx}{\linewidth}{@{}l c X c@{\hspace{1.4em}} l c X c@{}}
    \hline
    \textbf{(a)} $2^{x}=16$ & \blank{1.0cm} & \blank{1.8cm} & \blank{1.1cm} &
    \textbf{(b)} $2^{x}=12$ & \blank{1.0cm} & \blank{1.8cm} & \blank{1.1cm} \\
    \textbf{(c)} $5^{\,x+1}=\sqrt5$ & \blank{1.0cm} & \blank{1.8cm} & \blank{1.1cm} &
    \textbf{(d)} $7^{x}=1$ & \blank{1.0cm} & \blank{1.8cm} & \blank{1.1cm} \\
    \textbf{(e)} $10^{x}=45$ & \blank{1.0cm} & \blank{1.8cm} & \blank{1.1cm} &
    \textbf{(f)} $\left(\frac14\right)^{x}=8$ & \blank{1.0cm} & \blank{1.8cm} & \blank{1.1cm} \\
    \hline
    \end{tabularx}
    \end{center}
    You should have exactly two \textbf{no}s. For each of those two, answer this: does the equation
    still have a solution, and between which two whole numbers does it sit?
    \blank{7.0cm}

  \item \textbf{Error analysis.} Name the error, then give the correct solution.

    \textbf{(a)} Jamal solved $8^{x}=32$ this way: ``$\left(2^{3}\right)^{x}=2^{5}$, so
    $2^{\,3+x}=2^{5}$, so $3+x=5$ and $x=2$.'' \writelines{2}

    \textbf{(b)} Elise solved $9^{\,x+2}=27$ this way: ``$3^{\,2x+2}=3^{3}$, so $2x+2=3$ and
    $x=\frac12$.'' \writelines{2}

  \item \textbf{One answer or two?} Solve both equations, then explain the difference using the
    \emph{shape} of each graph --- not the algebra.

    \vspace{3pt}
    \hspace*{1.2em} $x^{2}=64$ \; gives $x=$ \blank{2.4cm} \hfill
    $2^{x}=64$ \; gives $x=$ \blank{2.4cm}

    \vspace{2pt}
    \writelines{2}
\end{enumerate}
\end{notesbox}

\begin{extensionbox}
\textbf{When matching bases tells you there is nothing to find.} Consider
\[
  4^{x}\cdot 2^{\,x+1} = 8^{\,x-1}.
\]
\begin{enumerate}[label=(\alph*), leftmargin=1.8em, itemsep=5pt]
  \item Write each side as a single power of $2$, then set the exponents equal. Something strange
        happens --- write down the equation you end up with and say what it tells you.
  \item Now rewrite each side as a \emph{multiple of} $8^{x}$. Fill in the two constants:
        the left side is \blank{1.4cm} $\cdot\,8^{x}$ and the right side is \blank{1.4cm}
        $\cdot\,8^{x}$.
  \item In Lesson 6.2's language, both sides are vertical stretches of the same curve $y=8^{x}$.
        Explain why two such curves can never intersect. (Hint: setting them equal leads to
        $\frac{15}{8}\cdot 8^{x}=0$. Why is that impossible?)
  \item Compare this equation with $2^{x}=5$. Both defeated today's method, but for completely
        different reasons. Say what each reason is, in one sentence apiece --- and say which of the
        two equations actually has a solution.
\end{enumerate}
\writelines{6}
\end{extensionbox}

\begin{spiralbox}
\textbf{Coming next --- Lesson 6.4: Compound Interest, Half-Life, and $e$.} Today's method is fast, and
it is also narrow: it works only when both sides happen to be powers of one base. Next lesson the
equations come from money and medicine, where nobody arranges the numbers for you --- an investment of
\$1000 at $5\%$ interest reaches \$2000 when $1000(1.05)^{t}=2000$, and $2$ is not a power of $1.05$.
You will build the models, evaluate them, watch compounding more and more often produce the number
$e$, and then run into that wall on purpose. Keep the list you started today: $2^{x}=5$, the coffee's
$132(0.85)^{t}=32$, the soda's $(0.9)^{t}=2.118$, and now $1000(1.05)^{t}=2000$. Every one of them has
an answer. Unit 7 is where you finally get to write it down.
\end{spiralbox}

\end{document}
